


Generated by \texttt{python scripts/vf2\_duty.py} (keys \texttt{vf2.duty.*}). Labels: ANALYTICAL for the formulas, SIMULATED for the check against the stored ns-3 runs. Nothing was rerun. Time on air is \texttt{crashguard.lora.airtime\_ms} (Semtech formula, 125 kHz, coding rate 4/5, explicit header, CRC, 8 preamble symbols).

\subsection*{1. Which duty model each part of the code enforces}

{\scriptsize\raggedright
\begin{longtable}{p{0.237\linewidth}p{0.237\linewidth}p{0.237\linewidth}p{0.237\linewidth}}
\toprule
Part & Duty model enforced & Where & Channels or sub-bands \\
\midrule
\endhead
ns-3 application (\texttt{xsim\_alert.cc}), policies p1, quick, d1, r3s, r3t, r3a, r6m and others & \textbf{hourly allowance}: the airtime of every uplink in any trailing 3,600 s must stay within \texttt{duty x 3600} s (36 s at 1\%); the application waits (\texttt{WaitUntil}) before a send; the end-of-run line \texttt{DUTY} reports the largest trailing-hour use & \texttt{WaitUntil}, \texttt{Used}, \texttt{g\_limit = duty x 3600} & the LoRaWAN module is configured with the ALOHA region: one sub-band (868.0 to 868.6 MHz) whose duty cycle is \textbf{1} and one logical channel (868.1 MHz). The stock formula next start = start + airtime / duty therefore adds no off-time (duty 1). RX2 downlinks outside the known sub-band are not duty-limited after the local patch (\texttt{XSIM\_BUILD.md}) \\
custom simulator (\texttt{world\_sim.Window}) & hourly allowance (trailing window) & \texttt{Window.ok}, \texttt{wait\_until} & one channel (same capture rule); no sub-band model \\
trace replay (\texttt{r3\_t1}, \texttt{r5\_l4}) & hourly allowance of the device (\texttt{LIMIT\_S = 0.01 x 3600}; counter \texttt{duty\_exceeded}), no per-transmission off-time & \texttt{r3\_t1.py} lines 38 and 115 & one link per device; no crowd \\
433 MHz & \textbf{no channel plan or sub-band list exists in the repository}; the 1\% limit is an assumption and the Indian 433 MHz rule (G.S.R. 347(E), rule 3 row 1) allows a duty cycle of at most 10\% with no channel plan (\texttt{FINALIZE\_INDIA\_RULES.md}) & none & not modelled: the rules give no channel plan, so the number of channels a node may use is not limited by them \\
\bottomrule
\end{longtable}}

\textbf{Consequence.} Every stored result of this study (R3, R6-1 and the others) was produced under an hourly allowance with no per-transmission off-time. The lockout trap does \textbf{not arise under the evaluated model}. It arises in regions that enforce an off-time after each frame on a shared sub-band (for example the European 1\% sub-bands, where T\_off = T\_on x (1/duty - 1) = 99 x T\_on at 1\%).

\subsection*{2. Exact numbers for a region with per-transmission off-time (not evaluated)}

Inputs: duty d = 1\% (and 10\% for comparison); the confirm is due at T\_F = 1.808 s after the trigger (1.8 s window plus 8 ms); the flash starts at s\_f = t\_E + 8 ms, with t\_E the decision time of the early-exit rule on the held-out crash events (median 50 ms, mean 112.4 ms; keys \texttt{r6.r2.held.median\_time\_fixed\_ms}, \texttt{mean\_time\_fixed\_ms}). One sub-band shared by all frames of the node.

{\scriptsize\raggedright
\begin{longtable}{p{0.106\linewidth}p{0.106\linewidth}p{0.106\linewidth}p{0.106\linewidth}p{0.106\linewidth}p{0.106\linewidth}p{0.106\linewidth}p{0.106\linewidth}p{0.106\linewidth}}
\toprule
SF & A14 (ms) & 99 x A14 (s) & A7 (ms) & 99 x A7 (s) & A97 (ms) & 99 x A97 (s) & earliest confirm after the flash, median t\_E (s) & delay beyond T\_F (s) \\
\midrule
\endhead
7 & 46.34 & 4.59 & 36.10 & 3.57 & 169.22 & 16.75 & 4.69 & 2.88 \\
8 & 82.43 & 8.16 & 72.19 & 7.15 & 297.47 & 29.45 & 8.30 & 6.49 \\
9 & 164.86 & 16.32 & 123.90 & 12.27 & 533.50 & 52.82 & 16.54 & 14.74 \\
\bottomrule
\end{longtable}}

Rows are duty 1\%. The 13-byte ns-3 micro frame has the airtime of the 14-byte flash at SF7 and SF9 (\texttt{FRAME\_LAYOUT.md}), so its lockout equals the flash's. The same table at 10\% (earliest confirm, median t\_E): SF7 0.52 s, SF8 0.88 s, SF9 1.71 s, SF10 2.95 s.

\subsection*{3. The condition (ANALYTICAL)}

\textbf{Terms.} d: duty fraction of the sub-band; A\_f: time on air of the provisional frame; s\_f: start time of the flash; T\_F: time at which the confirm is due; A\_c: time on air of the confirm frame (97 bytes, or the micro frame); L = s\_f + A\_f / d: the earliest start of any further frame on the sub-band after the flash (per-transmission off-time model above).

\textbf{Proposition (a condition, not a theorem).} Under that model, with one sub-band for all frames, the confirm starts at max(T\_F, L), so it reaches the base at max(T\_F, L) + A\_c when the initial attempt succeeds. Hence:
\begin{itemize}
\item \emph{Severity content (A\_c = 97 bytes).} The confirm is delayed by (L - T\_F)+ against the single-stage alert, and the delay is zero \textbf{if and only if} L <= T\_F, that is A\_f <= d (T\_F - s\_f). With a 97-byte confirm the severity never arrives earlier than in the single-stage alert: the flash buys only the earlier notice that a crash happened (earlier by T\_F + A\_97 - s\_f - A\_f).
\item \emph{Micro-confirm (A\_c = A\_micro < A\_97).} The micro scheme delivers the severity content earlier than the single-stage alert \textbf{if and only if} L - T\_F < A\_97 - A\_micro.
\item \emph{Net reading of the flash.} The flash gains G = T\_F + A\_97 - s\_f - A\_f on the crash notice and costs D = (L - T\_F)+ on the severity content. The lockout it causes (L - T\_F) is shorter than the time it gains when L - T\_F < G.
\end{itemize}

\emph{Derivation (every step).} (a) The model fixes the earliest next start on the sub-band as L (the stock ns-3 rule and the European off-time rule agree). (b) The confirm is ready at T\_F and no frame may start before L, so its start is max(T\_F, L): both bounds are necessary and together sufficient. (c) The arrival of the initial attempt is the start plus A\_c. (d) The difference to the single-stage arrival T\_F + A\_97 is (L - T\_F)+ + A\_c - A\_97, which gives 1 and 2 as equivalences (the inequalities are between the same quantities in both directions, with no other assumption). (e) Item 3 is the difference of two sums stated in the text. A retry of the flash restarts the lockout and only adds to L, so L is a lower bound when retries occur.

\textbf{Numbers (computed, not copied).}

{\scriptsize\raggedright
\begin{longtable}{p{0.119\linewidth}p{0.119\linewidth}p{0.119\linewidth}p{0.119\linewidth}p{0.119\linewidth}p{0.119\linewidth}p{0.119\linewidth}p{0.119\linewidth}}
\toprule
SF & d & A14 (s) & critical airtime d (T\_F - s\_f) (s) & condition A (median t\_E) & condition A (mean t\_E) & condition C (median t\_E) & attempts of the 97-byte alert before 10 s \\
\midrule
\endhead
7 & 1\% & 0.0463 & 0.0175 & fails & fails & fails & 1 \\
7 & 10\% & 0.0463 & 0.1750 & holds & holds & holds & 6 \\
8 & 1\% & 0.0824 & 0.0175 & fails & fails & fails & 1 \\
8 & 10\% & 0.0824 & 0.1750 & holds & holds & holds & 4 \\
9 & 1\% & 0.1649 & 0.0175 & fails & fails & fails & 1 \\
9 & 10\% & 0.1649 & 0.1750 & holds & holds & holds & 2 \\
10 & 1\% & 0.2888 & 0.0175 & fails & fails & fails & 1 \\
10 & 10\% & 0.2888 & 0.1750 & fails & fails & fails & 1 \\
\bottomrule
\end{longtable}}

At 1\% the critical airtime is 17.5 ms for the median decision time: no frame of 7 bytes or more fits at any spreading factor (A7 at SF7 is 36.1 ms), so with one sub-band and 1\% a flash always delays the confirm. At 10\% the condition holds for the 14-byte flash up to SF9 (median decision time) and up to SF9 (mean decision time).

\textbf{With more than one sub-band.} If the confirm is sent on a different sub-band from the flash, L does not apply to it and the delay is zero. The number of sub-bands a node may use at 433 MHz is not known from the repository (section 1), so no claim is made for that case; a node that is allowed two sub-bands should send the flash on one and the confirm on the other.

\subsection*{4. Check against the stored ns-3 results}

{\scriptsize\raggedright
\begin{longtable}{p{0.158\linewidth}p{0.158\linewidth}p{0.158\linewidth}p{0.158\linewidth}p{0.158\linewidth}p{0.158\linewidth}}
\toprule
SF & Policy & Two-stage events with flash and confirm delivered & Share whose confirm started inside the lockout & Median gap (confirm start - flash start) (s) & Gap required by the off-time model (100 x A14) (s) \\
\midrule
\endhead
7 & r3t & 114,968 & 1.0000 & 1.750 & 4.634 \\
7 & r6m & 115,996 & 1.0000 & 1.750 & 4.634 \\
9 & r3t & 108,055 & 1.0000 & 1.750 & 16.486 \\
9 & r6m & 112,286 & 1.0000 & 1.750 & 16.486 \\
\bottomrule
\end{longtable}}

\textbf{Reading.} The stored results are consistent with the proposition in its no-lockout branch, and only there: they were produced with a sub-band duty of 1 (ALOHA region of the module) and an hourly allowance, so the confirm starts at its due time T\_F whatever the flash did, and every one of the events above would have been delayed under a per-transmission off-time at 1\%. This is the reason, and nothing was rerun to make the two agree. The stored two-stage delivery and extra-airtime numbers therefore \textbf{do not transfer} to a region that enforces an off-time on one shared sub-band.

\subsection*{5. One link-aware gating rule (merged with the spreading-factor rule of R6-7)}

Enable the two-stage alert (flash, then confirm) at a spreading factor only if \textbf{both} hold:
\begin{itemize}
\item \emph{Load gate (R6-7, ROUND6\_BOUNDS.md section 2):} the lower bound of the extra airtime times the retry and collision inflation measured in ns-3 stays within the 40\% limit. Evaluated: SF7 passes, SF9 fails; SF8 and SF10 to SF12 were not run.
\item \emph{Lockout gate (this section), only in a region that enforces a per-transmission off-time on the sub-band that carries both frames:} condition A holds (A\_f <= d (T\_F - s\_f)), or the confirm is sent on another sub-band. Under an hourly allowance only, this gate is open at every spreading factor.
\end{itemize}

Rule: send two-stage alerts up to the highest spreading factor at which both gates are open, and the single-stage full decision above it. Under the evaluated model (hourly allowance) the rule gives SF7; under a 1\% per-transmission off-time on one sub-band it gives no spreading factor (single-stage everywhere); under 10\% it gives SF7 by the load gate alone. This is a design rule derived from the analysis and checked against two simulated cells and the stored runs above; it is not a tested hypothesis.

